\expandafter\gdef\csname SNside1\endcsname{\node[anchor=north west,inner sep=0pt,text width=48mm,font=\SideFont,align=left] at ([xshift=138mm,yshift=-500.62pt]current page.north west) {\textcolor{SNink}{\textbf{最小反例 · 空格仍是字符}}\par\vspace{1.5mm}空串长度为0，一个空格构成的串长度为1。仅凭打印后看起来空白无法判断串是否为空；接口判空应检查长度或表示规则，而不是观察屏幕上的可见字形。};}
\expandafter\gdef\csname SNside2\endcsname{\node[anchor=north west,inner sep=0pt,text width=48mm,font=\SideFont,align=left] at ([xshift=138mm,yshift=-252.67pt]current page.north west) {\textcolor{SNink}{\textbf{错项排除 · 一次候选失败才移起点}}\par\vspace{1.5mm}朴素法在某候选起点内逐字符比较，失配后起点才加1；主串和模式的工作指针可以回退。不能把朴素法的``重试下一个起点''与KMP保持i不动混为同一规则。};}
\expandafter\gdef\csname SNside3\endcsname{\node[anchor=north west,inner sep=0pt,text width=48mm,font=\SideFont,align=left] at ([xshift=138mm,yshift=-95.00pt]current page.north west) {\textcolor{SNink}{\textbf{不变量 · 回退保留相等边界}}\par\vspace{1.5mm}失配时已知T{[}i\ensuremath{-}j..i\ensuremath{-}1{]}等于P{[}0..j\ensuremath{-}1{]}。新候选只可保留旧段中同时是模式前缀的后缀；next{[}j{]}给出最长可保留长度，主串当前字符尚未被丢弃。};}
\expandafter\gdef\csname SNside4\endcsname{\node[anchor=north west,inner sep=0pt,text width=48mm,font=\SideFont,align=left] at ([xshift=138mm,yshift=-95.00pt]current page.north west) {\textcolor{SNink}{\textbf{错项排除 · next不含当前失配字符}}\par\vspace{1.5mm}本章next{[}j{]}研究P{[}0..j\ensuremath{-}1{]}的真前后缀，而不是P{[}0..j{]}。例中失配下标5之前是abaab，最长边界ab长2；把c算进去会让表错位。};}
\expandafter\gdef\csname SNside5\endcsname{\node[anchor=north west,inner sep=0pt,text width=48mm,font=\SideFont,align=left] at ([xshift=138mm,yshift=-95.00pt]current page.north west) {\textcolor{SNink}{\textbf{不变量 · 只跳过必然失败的位置}}\par\vspace{1.5mm}若P{[}next{[}j{]}{]}与已经失配的P{[}j{]}相同，那么同一主串字符与它再比较仍失败。nextval压缩这类失败链；只有获得这一字符相等条件，才能证明跳过不会漏解。};}
\expandafter\gdef\csname SNside6\endcsname{\node[anchor=north west,inner sep=0pt,text width=48mm,font=\SideFont,align=left] at ([xshift=138mm,yshift=-208.89pt]current page.north west) {\textcolor{SNink}{\textbf{最小反例 · 长串未必大}}\par\vspace{1.5mm}按通常英文编码的字典序，z大于aa，虽然前者只有一个字符。第一次不同字符就决定比较结果；只有短串完全等于长串前缀时，串长才参与最后判断。\par\smallskip\begin{tikzpicture}[x=1mm,y=1mm,>=Stealth,node distance=5mm]
\node[draw=SNink!50,fill=SNblue,rounded corners=1mm,text width=34mm,align=center,inner sep=1.2mm,font=\fontsize{7.5}{9}\selectfont] (m0) at (0,0) {z：长1};
\node[draw=SNink!50,fill=SNblue,rounded corners=1mm,text width=34mm,align=center,inner sep=1.2mm,font=\fontsize{7.5}{9}\selectfont] (m1) at (0,-11) {首字符z>a};
\node[draw=SNink!50,fill=SNblue,rounded corners=1mm,text width=34mm,align=center,inner sep=1.2mm,font=\fontsize{7.5}{9}\selectfont] (m2) at (0,-22) {aa：长2};
\draw[->,SNink] (m0) -- (m1);
\draw[->,SNink] (m1) -- (m2);
\end{tikzpicture}
\par{\fontsize{7.5}{9}\selectfont 字典序先看首个不同字符，z大于aa}};}
\expandafter\gdef\csname SNside7\endcsname{\node[anchor=north west,inner sep=0pt,text width=48mm,font=\SideFont,align=left] at ([xshift=138mm,yshift=-95.00pt]current page.north west) {\textcolor{SNink}{\textbf{课程连接 · 编码单位决定下标}}\par\vspace{1.5mm}字符编码把符号落实为存储。UTF-8的一个字符可能占多个字节，字节下标不总是字符位置；本章英文题按单字节模型分析，换输入模型时先确定元素单位。};}
\expandafter\gdef\csname SNside8\endcsname{\node[anchor=north west,inner sep=0pt,text width=48mm,font=\SideFont,align=left] at ([xshift=138mm,yshift=-95.00pt]current page.north west) {\textcolor{SNink}{\textbf{错项排除 · 不会变小才是准确表述}}\par\vspace{1.5mm}标准KMP中i在匹配成功或哨兵推进时变大，在模式回退时不动。因此``不会变大''错误，``可能变小''也错误；主串指针的单调不减与模式指针的失败回退必须分开看。};}
\expandafter\gdef\csname SNside9\endcsname{\node[anchor=north west,inner sep=0pt,text width=48mm,font=\SideFont,align=left] at ([xshift=138mm,yshift=-95.00pt]current page.north west) {\textcolor{SNink}{\textbf{不变量 · 半开区间分清空与越界}}\par\vspace{1.5mm}从pos取len个字符表示半开区间{[}pos,pos+len)。pos等于n时只可取长度0；检查len不超过n\ensuremath{-}pos，既表明边界，又避免先计算pos+len可能产生的整数溢出。};}
\AddToShipoutPictureFG{\begin{tikzpicture}[remember picture,overlay]\ifcsname SNside\thepage\endcsname\csname SNside\thepage\endcsname\fi\end{tikzpicture}}